\frontchapter{符号与缩略语}
\begingroup\small\setstretch{1.0}
\begin{longtable}{@{}>{\raggedright\arraybackslash}p{.25\linewidth}>{\raggedright\arraybackslash}p{.70\linewidth}@{}}
\toprule
符号 & 含义 \\
\midrule
$N,d$ & 数据库向量数与原始维度 \\
$M,D_s$ & 主码子空间数与每子空间维度 \\
$a_{\rm add},a_{\rm in},a_{\rm assign}$ & 配置 add 行数、输入暂存行数及分配评分缓冲的独立容量（行数） \\
$n_b,S_{\rm stage}$ & 流水缓冲数及每个输入／旋转数组的字节预算 \\
$\Pi$ & 方形正交变换；$q'=\Pi q$ \\
$B,B_r$ & 每主码子空间位数与每残差坐标位数 \\
$\rho(p),o_\ell$ & 物理位置 $p$ 的输入 ID，及 IVF 列表 $\ell$ 的起点 \\
$k$ & 返回近邻数 \\
$n_c,c$ & 路由候选数及接受精排的主码保留数 \\
$\alpha$ & 精排倍率，$c=\max\{k,\lceil\alpha k\rceil\}$ \\
$S,\epsilon$ & 样本数；允许缺失的参考槽位数 \\
$C_\alpha,\mathcal R_\alpha$ & 有效保留集合；预算 $\alpha$ 下返回的 ID \\
$\alpha_{\rm ref},B_{\rm reuse}$ & 参考精排预算与后续复用 batch 数 \\
$\glut$ & 分解查找表的组数 \\
$\gscan$ & 每查询扫描组数；零表示单 block 路径 \\
$C_{\rm buf},T_{\rm blk}$ & 缓冲容量；每 block 线程数 \\
$n_{\rm LUT}(\glut)$ & 每主码子空间 LUT 项数 \\
$\kappa(D,p)$ & 阈值使用的打包完整键 \\
$\lrefine$ & 每个保留候选的协作 lanes 数 \\
$\rawdist$ & $\|q-x\|_2^2$ \\
$\primarydist$ & $\|q'-\widehat p\|_2^2$ \\
$\refineddist$ & $\|q'-\widehat p-\widehat r\|_2^2+\|q'\|_2^2$ \\
\bottomrule
\end{longtable}

\bigskip
\begin{longtable}{@{}>{\raggedright\arraybackslash}p{.25\linewidth}>{\raggedright\arraybackslash}p{.70\linewidth}@{}}
\toprule
缩略语 & 含义 \\
\midrule
ANN & 近似最近邻 \\
ADC & 非对称距离计算 \\
IVF & 倒排文件 \\
JHQ & Johnson--Lindenstrauss 增强的层级量化 \\
LUT & 查找表 \\
PQ & 乘积量化 \\
QPS & 每秒查询数 \\
GT & 真实近邻 \\
FP32 & IEEE binary32；GEMM 模式另行说明 \\
TF32 & 低精度 Tensor Core 输入；FP32 累加 \\
CV & $\operatorname{sd}(n_c)/\operatorname{mean}(n_c)$ \\
\bottomrule
\end{longtable}
\endgroup
